\documentclass[11pt]{article}

% ---------------------------------------------------------------------------
%  Fixed polyplets through a(34): sixteen new terms of A006770, symmetry
%  companions, and the diagonal structure of king-animal triangles.
%  (Restructured 2026-07-05 per paper/restructure-plan.md: results first,
%  explicit confidence tiers, methods as an escalation ladder.)
% ---------------------------------------------------------------------------
\usepackage[margin=1in]{geometry}
\usepackage{amsmath,amssymb}
\usepackage{amsthm}
\usepackage{booktabs}
\usepackage{graphicx}
\usepackage{numprint}
\usepackage{microtype}
\usepackage[hidelinks]{hyperref}
\usepackage{flafter}   % a float never appears before the text that introduces it
\usepackage{float}     % [H] pins tables/figures exactly where written
\usepackage[font=small,labelfont=bf,labelsep=period]{caption}
\setlength{\intextsep}{16pt plus 3pt minus 3pt}

\npthousandsep{,}
\npfourdigitsep
\newcommand{\num}[1]{\numprint{#1}}
\newcommand{\byH}{\ensuremath{B_H}}
\newcommand{\tierone}{\textsf{T1}}
\newcommand{\tiertwo}{\textsf{T2}}
\newcommand{\tierthree}{\textsf{T3}}
\newcommand{\tiertwominus}{\textsf{T2}$^{-}$}
\newtheorem{conjecture}{Conjecture}
\newtheorem{openproblem}[conjecture]{Open Problem}

\title{\textbf{Fixed polyplets through $a(34)$: sixteen new terms of OEIS
A006770, their symmetry companions, and the diagonal structure of king-animal
triangles}\\[4pt]
\large A computational report}
\author{Jason Parker \texttt{<jasonp@panix.com>}\\[2pt]
\small Independent researcher (unaffiliated)}
\date{July 2026}

\begin{document}
\maketitle

\begin{abstract}
\noindent
A \emph{polyplet} (king-move animal) is a finite set of unit cells of the square
lattice connected through shared edges or corners. OEIS sequence
\href{https://oeis.org/A006770}{A006770} counts the \emph{fixed} polyplets of $n$
cells, previously known through $n=18$. We report sixteen new terms, through
\[
  a(34) = \num{515316838423862758858377704},
\]
each computed from scratch and reproducing every prior term. Because claims of
this kind are only as good as their verification, every number in this paper
carries an explicit confidence tier: $a(19)$ is confirmed by two algorithms
sharing no counting logic; $a(20)$--$a(33)$ by a single algorithm family
decorrelated across architectures and cross-checked per term against held-out
closed forms for the top diagonals of the height triangle (once certified by
a live held-out prediction, those forms also supply the tallest, sparsest
strata---$6.5\%$ of $a(34)$ arrives through them, a dependency quantified
exactly in the text); $a(34)$'s own top diagonal has no held-out closed
form, which we state rather than paper over. The one-sided companion A030233 is extended to
$n=34$; the free, bilateral, and asymmetric companions (A030222, A030234,
A030235, A194596) to the reach of the diagonal-mirror symmetry count. Along the
way the machinery yields structural results: closed forms for the diagonals
of the height triangle (an exponential/cumulant law with two new constants
per level---empirically pinned and validated on every held-out term, not
proved), a conjectural quasi-polynomial law for the diagonals of the
symmetric-fixed-point triangle (period two, degree $k$, with
generating-function numerators satisfying $N_k(\pm1)=(\pm2)^k$), a hole-count
stratification of A006770, exact fixed-height generating functions, the
extremal maximum hole area ($\sim n^2/8$), and rigorous lower bounds on
$a(n)$. Results computed and
validated but not proved are labeled as such; all code, data, and a
claim-verifier are released.
\end{abstract}

\section{Introduction}
\label{sec:intro}

The enumeration of lattice animals---finite connected sets of cells on a regular
lattice---is a classical problem at the meeting point of combinatorics and
statistical physics~\cite{klarner,jensen2001}. On the square lattice the two
standard families are \emph{polyominoes}, connected through shared edges, and
\emph{polyplets}, connected through shared edges or corners. The fixed counts
grow exponentially, so exact values are known only for modest $n$ and each new
term is a small computational milestone. The fixed-polyomino sequence
(\href{https://oeis.org/A001168}{A001168}) has been pushed
furthest~\cite{redelmeier,jensen2003,barequet}; the fixed-polyplet sequence
A006770 had been recorded through $n=18$.

This document is a \emph{computational report}: its claims are numbers, the
provenance and verification of those numbers, and the empirical structure
found in them. It reports the terms $a(19)$ through $a(34)$ and the results
that grew out of computing them, including two structural observations about
the diagonals of king-animal triangles (Section~\ref{sec:diagonals}), stated
as validated conjectures since proofs are not attempted here. It is an
amateur contribution carried out with substantial AI assistance
(Section~\ref{sec:ai}). The data appears in the document as well as the
repository: a report's job is to be usable by the next person to attempt
these computations.

\paragraph{Contributions.}
\begin{enumerate}
  \item Sixteen new terms of A006770, $a(19)$--$a(34)$
        (Section~\ref{sec:result}).
  \item The one-sided companion A030233 to $n=34$; the free/bilateral/
        asymmetric companions A030222, A030234, A030235, and A194596 toward
        $n=32$ (Section~\ref{sec:companions}).
        % TODO(n32-companions): the dmirror n=32 strip farm completes today;
        % insert the four companion values at n=25..32 when banked, and
        % update the reach statement. n=33 (T3, P_4-assisted) is running.
  \item Closed forms, empirically pinned and holdout-validated (not
        proved), for the height-triangle diagonals $T(n,n-k)$, $k\le15$,
        via an exponential/cumulant law, used both to accelerate the sweeps
        and---critically---as \emph{held-out per-term checks} of the
        frontier terms (Sections~\ref{sec:diagonals}, \ref{sec:confirm}).
  \item A conjectural quasi-polynomial law for the diagonals of the diagonal-mirror
        symmetry triangle, with pinned polynomials, forward confirmations,
        and generating-function structure (Section~\ref{sec:diagonals}).
  \item A hole-count stratification of A006770 exact through $n=18$, the
        extremal hole area $M(n)\sim n^2/8$, exact fixed-height generating
        functions for $H\le10$, rigorous lower bounds, and a uniform-sampling
        portrait (Sections~\ref{sec:holes}--\ref{sec:sampling}).
\end{enumerate}

\paragraph{Confidence tiers.}
Since a single long count from one program is properly a conjecture, every
numeric claim below is tagged with one of three tiers, and the results tables
repeat the tags. The tiers are honest about an asymmetry: past Redelmeier's
reach no second counting \emph{algorithm} for the full sequence exists, so the
later terms rest on decorrelation rather than algorithmic independence.

\begin{table}[H]
  \centering\small
  \begin{tabular}{l p{0.72\textwidth}}
    \toprule
    tier & meaning \\
    \midrule
    \tierone{} & Two algorithms sharing no counting logic agree, each
    reproducing the full published range. Covers $a(1)$--$a(19)$; the
    symmetry counts to $n=24$; everything derived only from these. \\
    \tiertwo{} & One counting algorithm, decorrelated and cross-checked:
    full chain-reproduction of all prior terms on every run; independent
    recount on a second architecture/compiler; agreement of the top height
    strata with \emph{held-out} diagonal closed forms $P_k$ (fitted on
    earlier terms, predicting cells of later ones); modular/CRT internal
    consistency. Covers $a(20)$--$a(33)$ and the symmetry counts
    $25\le n\le32$. \\
    \tiertwominus{} & \tiertwo{} \emph{minus} the held-out diagonal check.
    Covers exactly one number, $a(34)$: its top surviving stratum would be
    checked by $P_{16}$, which has no independent holdout until $a(35)$
    exists. It gets its own tag so the weaker standard is never silently
    absorbed into \tiertwo{}. \\
    \tierthree{} & Exact computation composed with \emph{empirically pinned
    but unproven} quasi-polynomial formulas (the diagonal-mirror $P_k$ of
    Section~\ref{sec:diagonals}, each pinned with holdout hits and confirmed
    forward on strips computed after pinning). The failure mode is specific
    and stated: an onset shift or degree transition beyond the pinned range
    would silently corrupt exactly the substituted strips; the forward
    confirmations bound this risk but cannot eliminate it. Planned for the
    $n=33$ companions; no \tierthree{} value is presented as established in
    this paper without its label. \\
    \bottomrule
  \end{tabular}
  \caption{Confidence tiers used throughout. Conservative reading: \tierone{}
  is confirmed; \tiertwo{} is a decorrelated single-method result;
  \tierthree{} is conjecture-assisted and labeled as such.}
  \label{tab:tiers}
\end{table}

\section{Definitions}
\label{sec:object}

Fix the square lattice $\mathbb{Z}^2$ and call its unit squares \emph{cells}. A
\emph{lattice animal} is a finite, nonempty, adjacency-connected set of cells.
Under \textbf{rook} (edge) adjacency the animals are \emph{polyominoes}; under
\textbf{king} (edge-or-corner) adjacency they are \emph{polyplets}. Every
polyomino is a polyplet but not conversely: two cells touching only at a corner
form a polyplet and not a polyomino. One subtlety: a
polyplet connected \emph{purely} by diagonals lives on a single checkerboard
colour class, itself a square lattice, so such a shape is a polyomino in
disguise. A genuinely king shape mixes edge and corner contacts.

Animals are counted \textbf{fixed} (up to translation; rotations and
reflections distinct), \textbf{one-sided} (up to translation and rotation), or
\textbf{free} (up to all isometries). A006770 is the fixed count. Two
triangles refine it: the \emph{height triangle} $T(n,H)$ (animals of $n$ cells
and bounding-box height exactly $H$, so $a(n)=\sum_H T(n,H)$), and, for each
symmetry $g$ of the square, the count of fixed animals invariant under $g$,
which enters the companion counts by Burnside's lemma
(Section~\ref{sec:companions}).

\section{Headline results}
\label{sec:result}

\begin{equation}
  \boxed{\,a(34) = \num{515316838423862758858377704}\,}
\end{equation}

Table~\ref{tab:terms} gives the full sequence. Terms $n\le18$ reproduce
A006770 exactly; $a(19)$--$a(34)$ are new. $a(19)$ is \tierone{};
$a(20)$--$a(33)$ are \tiertwo{}; $a(34)$ is \tiertwominus{}
(Table~\ref{tab:tiers}), the most exposed number in this paper.

\begin{table}[H]
  \centering
  \small
  \begin{tabular}{r r @{\quad}|@{\quad} r r}
    \toprule
    $n$ & $a(n)$ & $n$ & $a(n)$\\
    \midrule
    1  & \num{1}              & 18 & \num{22471158811164}\\
    2  & \num{4}              & \textbf{19} & $\mathbf{\num{151609203011580}}$\\
    3  & \num{20}             & \textbf{20} & $\mathbf{\num{1025573519362016}}$\\
    4  & \num{110}            & \textbf{21} & $\mathbf{\num{6954084405510437}}$\\
    5  & \num{638}            & \textbf{22} & $\mathbf{\num{47255332844367680}}$\\
    6  & \num{3832}           & \textbf{23} & $\mathbf{\num{321749260511448732}}$\\
    7  & \num{23592}          & \textbf{24} & $\mathbf{\num{2194666793369310473}}$\\
    8  & \num{147941}         & \textbf{25} & $\mathbf{\num{14994811325186658577}}$\\
    9  & \num{940982}         & \textbf{26} & $\mathbf{\num{102607513847014153892}}$\\
    10 & \num{6053180}        & \textbf{27} & $\mathbf{\num{703126792093436034256}}$\\
    11 & \num{39299408}       & \textbf{28} & $\mathbf{\num{4824589228356722264087}}$\\
    12 & \num{257105146}      & \textbf{29} & $\mathbf{\num{33145129805782422061325}}$\\
    13 & \num{1692931066}     & \textbf{30} & $\mathbf{\num{227969227118066423789154}}$\\
    14 & \num{11208974860}    & \textbf{31} & $\mathbf{\num{1569631619081324581014300}}$\\
    15 & \num{74570549714}    & \textbf{32} & $\mathbf{\num{10818203977457804418974036}}$\\
    16 & \num{498174818986}   & \textbf{33} & $\mathbf{\num{74631481980411777590683952}}$\\
    17 & \num{3340366308393}  & \textbf{34} & $\mathbf{\num{515316838423862758858377704}}$\\
    \bottomrule
  \end{tabular}
  \caption{Fixed polyplets $a(n)$, $n=1,\dots,34$. Terms $1$--$18$ match the
  published A006770; boldface terms are new. Tiers: $a(19)$ \tierone{};
  $a(20)$--$a(33)$ \tiertwo{}; $a(34)$ \tiertwominus{}.}
  \label{tab:terms}
\end{table}

\paragraph{Growth.}
The successive ratios $r_n=a(n)/a(n{-}1)$ climb slowly:
$r_{20}\approx6.765$, $r_{27}\approx6.853$, $r_{34}\approx6.905$. Fitting
the standard ansatz $r_n=\lambda(1+\theta/n)$ by least squares over
trailing windows gives estimates that are stable in the third decimal:
$\lambda=7.101$, $7.104$, $7.105$, $7.106$ (with $\theta$ between $-0.94$
and $-0.96$) for windows starting at $n=10$, $15$, $20$, $25$ and ending at
$34$. The
slight upward drift with window start is the visible confluent correction;
the exponent is consistent with the universal two-dimensional lattice-animal
value $\theta=-1$~\cite{jensen2001}. We accordingly estimate
$\lambda\approx7.10$--$7.11$, and emphasise that this is a fit without error
control---no differential-approximant analysis has been attempted---so it is
an estimate, not a bound.

\paragraph{One-sided companion.}
The one-sided count (A030233) needs, besides the fixed count, only the
$90^\circ$- and $180^\circ$-rotation fixed-point counts
(Section~\ref{sec:companions}), both of which reach $n=34$; hence fifteen new
terms, through
\[
  \mathrm{OneSided}(34)=\num{128829209605977867230343869}.
\]
Table~\ref{tab:onesided} lists them. The remaining companions also need the
mirror counts and are bounded by the diagonal-mirror count's reach
(Section~\ref{sec:companions}).

\begin{table}[H]
  \centering
  \small
  \begin{tabular}{r r @{\quad}|@{\quad} r r}
    \toprule
    $n$ & A030233$(n)$ & $n$ & A030233$(n)$\\
    \midrule
    20 & \num{256393397108358}      & 28 & \num{1206147307126536440324}\\
    21 & \num{1738521118535082}     & 29 & \num{8286282451482505589564}\\
    22 & \num{11813833328245848}    & 30 & \num{56992306779773178547095}\\
    23 & \num{80437315244081648}    & 31 & \num{392407904770584270738465}\\
    24 & \num{548666699140240976}   & 32 & \num{2704550994366217058008641}\\
    25 & \num{3748702832087014098}  & 33 & \num{18657870495104684580672401}\\
    26 & \num{25651878467205504568} & 34 & \num{128829209605977867230343869}\\
    27 & \num{175781698028751634291}& & \\
    \bottomrule
  \end{tabular}
  \caption{New terms of the one-sided count A030233 (previously known to
  $n=19$ via this project's earlier terms; to $n=17$ before that). Tier:
  \tierone{} through $n=24$, \tiertwo{} beyond ($n=34$ inherits the fixed
  count's $a(34)$ caveat). Every derivation validates the $/4$ integrality of
  the Burnside numerator and every published OEIS term.}
  \label{tab:onesided}
\end{table}

\section{Methods: an escalation ladder}
\label{sec:methods}

Four methods computed these numbers; each subsection states what ran out and
what replaced it. All engines share three disciplines: exact integer
arithmetic; every production run re-derives \emph{all} smaller terms, so a
regression anywhere in range is visible; and admissible pruning only (a prune
may only ever \emph{lose} animals, so reproducing the published range
certifies that it lost none).

\subsection{Method A: direct enumeration (to $n=19$)}
\label{sec:methodA}

Redelmeier's backtracking enumeration~\cite{redelmeier}, adapted to king
adjacency, grows every animal exactly once from an anchored cell: it maintains
an ordered frontier of \emph{untried} neighbours and recurses, with the
standard rule that a cell once placed or rejected at a level is excluded from
deeper levels. The recursion tallies by size; no animal is stored. Being
embarrassingly parallel by sub-tree, the count was run as two decorrelated
campaigns---on Apple silicon (Apple~clang, frontier split $7\times64$) and
AMD/x86 (GCC, split $8\times96$). The per-size outputs are
\textbf{byte-identical}. Because the two campaigns differ in hardware,
compiler, and scheduling and the arithmetic is exact integer addition, a
transient or platform-specific fault in one is overwhelmingly unlikely to be
mirrored digit-for-digit in the other; the error mode they cannot jointly
catch is a bug in the shared \emph{algorithm}, which is what Method~B rules
out at $n\le19$.

\subsection{Method B: the column transfer matrix (to $n\approx23$)}
\label{sec:methodB}

The second method counts the same animals by an entirely different principle:
dynamic programming over the connectivity of a moving boundary, with no animal
materialised---a king adaptation of the transfer-matrix / finite-lattice
technique for polyominoes~\cite{jensen2001,jensen2003,bousquet}. Animals of
bounding-box height exactly $H$ are counted separately, so
$a(n)=\sum_{H=1}^{n}T(n,H)$; the heights are independent, checkpointable
sub-problems.

Within a height, confine the animal to an $H$-row strip and sweep columns left
to right. A \emph{boundary signature}---the occupancy of the most recent
column, a partition of its filled cells into the partial animal's connected
components, and two flags for whether rows $0$ and $H-1$ have yet been
touched---is a Markov state, since a future cell attaches only to the current
column or later. Because king adjacency reaches the diagonal neighbours
$(c{-}1,r{\pm}1)$, the transfer moves an \emph{entire column} at once, keeping
the previous column wholly present; connectivity across the step is resolved
by union--find. Three invariants make accepted sweeps biject with fixed
polyplets: the first column is column $0$ (no horizontal double-counting); an
animal is tallied only with both touch-flags set (height exactly $H$); and
only when its boundary is a single component---partial states that can never
reconnect are pruned. A size-budget prune drops any state whose cells-so-far
plus an \emph{admissible} lower bound on cells-still-needed exceeds $n$.

Concretely, a state is a fixed-width byte string: one byte per strip row
holding the cell's connected-component label ($0$ for empty), with labels
\emph{canonically renumbered in first-occurrence order} after every
transition---the normalization that prevents equivalent partitions from
proliferating as distinct states---plus the two touch flags; the value
attached to a state is its vector of counts by size. States live in a flat
open-addressing hash map keyed by that string (about $300$ bytes per state
including counts at production sizes). The full layout is in the release's
\texttt{core/signature.h}.

Method B confirmed $a(19)$ against Method A (byte-identical totals) and
carried the sequence into the early twenties before its state space---which
keeps a whole column of connectivity alive per step---outgrew memory.

\subsection{Method C: the cell-at-a-time transfer matrix (to $n=34$)}
\label{sec:methodC}

The production engine for $a(24)$--$a(34)$ replaces the column step with a
\emph{cell-at-a-time} step: the boundary advances one cell per transition,
so the frontier is a staircase---$r$ cells of the new column above $H-r$
cells of the old---plus one retained north-west \emph{carry} cell, kept
because king adjacency lets the next cell reach diagonally back across the
kink in the staircase; without it the step would sever connectivity
information the lattice still provides. The state encoding is as in
Method~B (canonical component-label bytes over the staircase cells and the
carry, touch flags, counts), and each transition decides a single cell,
resolving connectivity by union--find over its at-most-three placed
neighbours. The finer step multiplies the number of transitions but shrinks
the live state space and, decisively, the volume that must be reshuffled
between passes (measured speedups over Method B grow from $\sim$29$\times$
at $H=12$ to $\sim$200$\times$ at $H=16$). At frontier sizes beyond RAM the
state map spills to disk: the frontier is hash-sharded, expanded shard by
shard, and merged by external passes, so the binding resource is disk
bandwidth rather than memory. Heights remain independent jobs; per-height
results are banked with their provenance and re-summed by a combiner that
refuses partial or duplicated inputs.

\paragraph{I/O integrity.}
A multi-hour external shuffle is itself a failure surface, so the disk
layer is engineered to fail closed rather than lose states silently. Every
spill file carries a self-describing header (magic, format version, byte
order, key width, and the \emph{record count}, so truncation---including
the short writes produced by resource exhaustion---is detected on read
rather than absorbed); readers and the merge refuse malformed or
count-mismatched inputs; the per-height combiner refuses missing, empty,
or duplicated shards and takes a lock against concurrent runs on the same
directory; and kept artifacts are listed in a manifest with per-file
checksums. These refusal paths are not assumed to work: the release's gate
suite fault-injects each failure mode (truncated file, missing shard,
duplicated rows, held lock) and requires the pipeline to abort loudly.

Each production run reproduces the complete triangle---every $T(n,H)$ for all
smaller $n$---so each new term arrives with a full-chain regression against
everything already known, including the \tierone{} range.

\paragraph{Which cells are computed and which are injected.}
This must be stated exactly, because the headline totals contain both. For
$a(34)$, heights $3\le H\le18$ were swept for real ($93.5\%$ of the total);
$H=1,2$ are elementary closed forms; and $H\ge19$---the diagonals
$k\le15$---were \emph{injected} from the fitted closed forms of
Section~\ref{sec:heightdiag}, contributing $6.5\%$ of $a(34)$. The
analogous split holds for the preceding terms, the injection boundary
advancing as each new $P_k$ is certified (at $a(30)$: real sweep
$H\!=\!3$--$17$, closed forms elsewhere). The
discipline governing injection is strict and uniform: a $P_k$ is fitted on
real-swept cells of earlier terms, must reproduce every other real-swept
cell on its diagonal (holdout), must then predict the \emph{next} term's
real-swept cell on that diagonal exactly---the newest diagonal is always
swept for real once more before its formula is trusted---and only after
surviving that live prediction is it injected for subsequent terms. An
injected formula failing at larger $n$ would therefore have to pass a
genuine forward prediction and then break one term later without warning;
possible, unproven-formula risk being what it is, but the exposure is
quantified ($6.5\%$ of the frontier term) and every injected formula
carries at least one live held-out success. Readers who want totals free
of any formula assistance may sum Table~\ref{tab:byheight34}'s rows
$H\le18$ against the banked per-height data; the \tierone{}/\tiertwo{}
range $n\le19$ contains no injected cells at all.

Table~\ref{tab:byheight34} gives the height decomposition of the frontier
term with the computed/injected boundary as stated above.

\begin{table}[H]
  \centering
  \small
  \begin{tabular}{r r @{\quad}|@{\quad} r r}
    \toprule
    $H$ & $T(34,H)$ & $H$ & $T(34,H)$\\
    \midrule
    1  & \num{1}                          & 18 & \num{23288787870043631158670332}\\
    2  & \num{8822750406819}              & 19 & \num{15111742807653801090985593}\\
    3  & \num{1450547121316690736}        & 20 & \num{9061341124316405172057950}\\
    4  & \num{997738306231538048056}      & 21 & \num{5002581174202331698460002}\\
    5  & \num{54325938673309776005354}    & 22 & \num{2531214280334205695134436}\\
    6  & \num{718450867038312966264504}   & 23 & \num{1167265695441145358152351}\\
    7  & \num{3938007987492904266811062}  & 24 & \num{487286890123087057884426}\\
    8  & \num{12136378135639937908810168} & 25 & \num{182652546462335637272478}\\
    9  & \num{25406165739609109253893844} & 26 & \num{60863544290648528714052}\\
    10 & \num{40756412152266215827765016} & 27 & \num{17807197708330934888844}\\
    11 & \num{54159404667976998554266361} & 28 & \num{4502990884375275679854}\\
    12 & \num{62670808373359868559352846} & 29 & \num{964076818028408805330}\\
    13 & \num{65125540284467338650988595} & 30 & \num{169895606813568532170}\\
    14 & \num{61866157347759886121825501} & 31 & \num{23657686681679437077}\\
    15 & \num{54235518293248262678336909} & 32 & \num{2440297953560629818}\\
    16 & \num{44075625952411715905525520} & 33 & \num{165742361336192445}\\
    17 & \num{33255837131885526055333731} & 34 & \num{5559060566555523}\\
    \midrule
    \multicolumn{4}{r}{$\sum_H T(34,H) = \num{515316838423862758858377704}$}\\
    \bottomrule
  \end{tabular}
  \caption{The height decomposition of $a(34)$, peaking at $H=13$. The
  extreme diagonal is the closed form $T(H,H)=3^{H-1}$
  ($T(34,34)=3^{33}$).}
  \label{tab:byheight34}
\end{table}

\subsection{Method D: symmetric transfer matrices (the companions)}
\label{sec:methodD}

The companion counts need, for each symmetry $g$ of the square, the number of
fixed animals invariant under $g$. An animal fixed by $g$ is determined by a
fundamental-domain slice of itself, so these counts behave like
$\sqrt{a(n)}$-scale problems and can reach $n=34$ where the full count cannot
be re-derived independently. Two engines compute them:

\begin{itemize}
  \item An \emph{orbit-graph Redelmeier} enumerator (explicit, per-animal):
        growth over cell-orbits under $\langle g\rangle$. Being explicit it is
        the cross-algorithm oracle, and it carried all four symmetry types to
        $n=24$.
  \item A family of \emph{symmetric transfer matrices}, one per conjugacy
        class: a palindromic-column sweep for the axis mirror; a half-plane
        sweep with a self-gluing seam for the $180^\circ$ rotation (with a
        transpose restriction, counting each animal from its wide
        orientation, that collapses the tall strips); and, for the diagonal
        mirror, an exact-bounding-box \emph{hook} sweep---the frontier is the
        L-shaped shell of a growing square, swept corner outward, with the
        mirror symmetry enforced by construction. The $90^\circ$ count, whose
        invariant animals have $\sim n/4$ free cells, stays within the
        explicit enumerator's reach through $n=34$.
\end{itemize}

The two engines agree---live, in the release's test gate---for every type at
small $n$, and byte-identical at the $n=24$ oracle; beyond $n=24$ the
transfer-matrix values are single-algorithm (\tiertwo{}), validated by prefix
against the oracle range on every run. The diagonal-mirror sweep is the
expensive one (its bounding box is forced square, so the frontier compresses
least); its production form stores the frontier in flat sharded arrays with
ranged count vectors, and farms disjoint exact-bounding-box strips across
machines, the partial totals summing exactly by construction.
% TODO(n32-companions): dmirror n=32 completes today (ayr tail); n=33
% (T3-assisted, strips S>=29 from P_4) is running on dalby. Insert final
% reach + values in Section~\ref{sec:companions} when banked.

\subsection{Method E: diagonal closed forms}
\label{sec:methodE}

The top diagonals of the height triangle---$T(n,n-k)$ for small $k$, the
tallest, sparsest strips---are the most expensive cells of a sweep and the
least populated. They turn out to have exact closed forms
(Section~\ref{sec:diagonals}): for each $k\le15$ a polynomial law, fitted from
two data points and then \emph{validated on held-out terms}, replaces the top
real sweep height. Beyond the speedup, these closed forms are per-term
independent checks of the production engine: a polynomial pinned on terms
$\le n_0$ that predicts a frontier cell of term $n_0{+}1$ exactly is evidence
no single-engine recount can provide. The same idea applied to the
diagonal-mirror triangle yields the quasi-polynomials that will assist the
$n=33$ companions (\tierthree{} when used).

\section{Validation architecture}
\label{sec:confirm}

The verification burden shifts across the sequence, and the checks are
designed around the weakest link at each stage.

\begin{itemize}
  \item \textbf{External truth.} Every method reproduces A006770 exactly for
        every $n\le18$, on every production run (the triangle re-derives all
        smaller terms as a by-product).
  \item \textbf{Method A $=$ Method B at $a(19)$.} Byte-identical totals from
        algorithmically disjoint counters---the \tierone{} anchor.
  \item \textbf{Cross-architecture recounts.} $a(34)$'s full triangle was
        recomputed on a second architecture and compiler (x86/GCC against
        ARM/clang production) and is byte-identical, row by row. Earlier
        terms had per-term recounts of the heaviest strata.
  \item \textbf{Held-out diagonal closed forms.} For each $k\le15$, $P_k$
        was fitted on the smallest admissible terms and then required to
        reproduce every later term's cell $T(n,n-k)$ exactly. The flagship
        instance: $P_{15}$, pinned before $a(33)$ existed, predicted
        $T(33,18)$ exactly. $a(34)$'s top cell would need $P_{16}$, which
        can be fitted but has no independent holdout until $a(35)$; we
        therefore do not count it as a check, and $a(34)$ carries the weaker
        tier \tiertwominus{}. Two mitigations are worth stating precisely:
        the diagonals $k\le15$ of $a(34)$ itself \emph{are} $P_k$-checked,
        and $a(34)$'s top surviving stratum runs the same code path that
        $P_{15}$ certified at $a(33)$---what lacks certification is its
        execution at the larger size, not the path.
  \item \textbf{Modular consistency.} Production sweeps run with residue
        counters and CRT reconstruction; independent residue systems must
        agree, catching overflow and transient corruption (not shared
        enumeration bugs, which the items above target).
  \item \textbf{Structural cross-checks.} The transfer matrices'
        (height,\,size) and (size,\,perimeter) marginals match the explicit
        enumeration's; rook- and bishop-connected sub-counts reproduce
        A001168; the $H$-cell height-$H$ diagonal is exactly $3^{H-1}$;
        Burnside numerators satisfy their $/4$ and $/8$ integrality at every
        term.
  \item \textbf{Fault injection.} Planted errors are caught (a dropped
        diagonal at $n=2$, a doubled count at $n=1$, a prune off-by-one at
        $n=10$) while a provably count-preserving change survives---the
        checks fail on wrong programs and pass on right ones.
  \item \textbf{Hygiene.} The engines are AddressSanitizer/UBSan clean; the
        multithreaded paths are bit-identical to serial; every kept result
        traces to a clean-tree commit stamp in a provenance ledger, and the
        cross-checks are encoded as gates that run before release.
\end{itemize}

\section{The diagonal structure of king-animal triangles}
\label{sec:diagonals}

Both triangles in this paper---the height triangle of the full count and the
exact-bounding-box triangle of the diagonal-mirror count---turn out to have
exactly solvable top diagonals. We record the laws, the fitting discipline,
and what is proved versus pinned.

\subsection{The height triangle: an exponential/cumulant law}
\label{sec:heightdiag}

A height-$H$ animal of exactly $H$ cells is one cell per row drifting by
$\{-1,0,+1\}$: $T(H,H)=3^{H-1}$. On the diagonal $k=n-H$ ($k$ cells above
minimum), the normalized counts
$P_k(n) = T(n,\,n-k)\,\cdot 3^{\,1+3k-n}$
satisfy
\[
  P_k(n) \;=\; [y^k]\,\exp\!\Big(\sum_{j\ge1}(a_j+b_j\,n)\,y^j\Big),
\]
an exponential of \emph{cumulants linear in $n$}: $b_j$ is the per-site
density of connected weight-$j$ defect clusters against the random-walk
background, $a_j$ the boundary correction---so each level costs exactly two
new constants, fitted from two data points and validated on all others. The
leading coefficient is $25^k/k!$ (twenty-five single-defect species per walk
site, an integer the fit must and does reproduce at every level), and $P_k$
times $k!$ must be integral. Levels $k\le15$ are fitted, validated, and wired
into production (Section~\ref{sec:methodE}); each was pinned on the earliest
admissible terms and holds on every later one. The cumulant form is standard
polymer-gas combinatorics; we have not proved the linearity of the cumulants
for this family, and treat the law as empirically exact within its validated
range.

\subsection{The diagonal-mirror triangle: quasi-polynomials, period two}
\label{sec:dmdiag}

Let $d(S,n)$ count the diagonal-mirror-fixed polyplets with $n$ cells and
bounding box exactly $S\times S$ (the box is forced square). The minimal such
animal is a bare diagonal, $n=S$, and there are exactly two ground states
(main and anti-diagonal). On the diagonal $k=n-S$ the data show, and exact
difference tables confirm over every available point:
\[
  d(S,\,S+k)\ \text{is quasi-polynomial in } S \text{ of period two and
  degree } k \text{ per parity class},
\]
with leading coefficient $S^k/k!$ and empirical onset $S\ge 2k+2$ (even) /
$2k+3$ (odd). The pinned polynomials for $k\le3$
(Table~\ref{tab:dmpk}), with $4$--$12$ holdout hits each, were subsequently
confirmed \emph{forward} on strips computed hours later on different hardware
($P_1$ hit $d(31,32)=38$; $P_2$ hit $d(30,32)=672$; $P_3$ hit
$d(29,32)=7456$); $P_4$ (even class) is pinned by two independent methods
that agree exactly.

\begin{table}[H]
  \centering\small
  \renewcommand{\arraystretch}{1.5}
  \begin{tabular}{c l l}
    \toprule
    $k$ & even $S$ & odd $S$\\
    \midrule
    0 & $2$ & $2$\\
    1 & $S+6$ & $S+7$\\
    2 & $\tfrac12S^2+7S+12$ & $\tfrac12S^2+6S+\tfrac{27}{2}$\\
    3 & $\tfrac16S^3+3S^2+\tfrac{40}{3}S+50$ &
        $\tfrac16S^3+\tfrac72S^2+\tfrac{83}{6}S+\tfrac{93}{2}$\\
    4 & $\tfrac1{24}S^4+\tfrac32S^3+\tfrac{28}{3}S^2+33S+180$ &
        $\tfrac1{24}S^4+\tfrac43S^3+\tfrac{97}{12}S^2+\tfrac{119}{3}S+\tfrac{1367}{8}$\,$^\dagger$\\
    5 & $\tfrac1{120}S^5+\tfrac5{12}S^4+4S^3+\tfrac{55}{3}S^2+\tfrac{2278}{15}S+570$\,$^\dagger$ &
        $\tfrac1{120}S^5+\tfrac{11}{24}S^4+\tfrac{19}{4}S^3+\tfrac{185}{12}S^2+\tfrac{18869}{120}S+\tfrac{4545}{8}$\,$^\dagger$\\
    \bottomrule
  \end{tabular}
  \caption{Diagonal-mirror diagonal polynomials $d(S,S+k)$, valid from the
  stated onsets. Every polynomial was pinned on the highest-$S$ points then
  required to reproduce all earlier in-regime points (holdout); $k\le3$
  additionally hit strips computed after pinning (forward confirmation), and
  $P_4$ (even) is pinned by two independent methods that agree exactly.
  Entries marked $^\dagger$ are fitted through the cumulant form of the
  repository's exp-fitter with fewer exact witnesses ($9$ at level~4, $4$ at
  level~5) and should be treated as one grade weaker than the pinned rows.}
  \label{tab:dmpk}
\end{table}
% TODO(pk-status): P_4 odd pins when the ayr S=23 strip lands (in flight);
% P_5 both parities are fitted via the cumulant form with 4 exact witnesses.
% Update this paragraph to the final pin/confirmation status before submit.

In the generating-function basis the structure is crisper: with
$G_k(x)=\sum_S d(S,S+k)\,x^S$,
\[
  G_k(x)\;=\;\frac{N_k(x)}{(1-x)^{k+1}(1+x)^{k}},\qquad
  N_k \in \mathbb{Z}[x],\ \deg N_k = 2k,\qquad
  N_k(\pm1)=(\pm2)^k .
\]
The integer numerators, for the record:
\begin{align*}
  N_0 &= 2\\
  N_1 &= 6+2x-6x^2\\
  N_2 &= 12+8x-16x^2-8x^3+8x^4\\
  N_3 &= 50+14x-124x^2-8x^3+110x^4+2x^5-36x^6\\
  N_4 &= 180+40x-644x^2-54x^3+942x^4+12x^5-600x^6+2x^7+138x^8\\
  N_5 &= 570+176x-2610x^2-520x^3+4996x^4+764x^5-4816x^6-496x^7\\
      &\qquad +2282x^8+108x^9-422x^{10}
\end{align*}
($N_4$ and $N_5$ are fitted rather than pinned, like their rows in
Table~\ref{tab:dmpk}).
The boundary values say the dominant single-defect family has two species per
odd footprint; the next layer of the peeled numerators follows linear-in-$k$
laws, the signature of interacting defect pairs. Exact overdetermined
elimination \emph{refutes} any low-order bivariate rational closure across
$k$ (each level carries genuinely new defect species), and the parity
structure of the $k=3$ level refutes the naive two-ground-state defect gas:
additional length-free families (anti-diagonal excursions crossing the main
spine) are present. A proof is not attempted, and we are explicit about
what one would require. The natural strategy is a decomposition argument:
show that every sufficiently large animal on this diagonal splits, in
exactly one way, into a bounded set of local pieces (stretches of clean
diagonal, the excursions above, and small clusters carrying the $k$ extra
cells), so that counting animals reduces to counting arrangements of pieces
along a line. Decompositions of that kind yield rational generating
functions; but rationality alone does not give quasi-polynomiality---for
that, every pole of $G_k$ must lie at a root of unity, and the specific law
observed here demands the poles sit at $x=\pm1$ with multiplicities exactly
$k{+}1$ and $k$. Nothing in the decomposition picture forces those
multiplicities; they are measured facts a proof would have to reproduce,
presumably from the count of independent free lengths available at weight
$k$. So the strategy is recorded as a strategy: the honest state of
knowledge is the validated data above, including the parity anomaly that
already ruled out the naive version of the argument.

All of Section~\ref{sec:dmdiag} is computed and validated, not proved. Its
practical role is sharp: the sparse strips ($S\ge29$ at $n\ge33$) are
precisely the RAM-expensive ones---one strip burned $\sim$$62$ CPU-hours and
$79$\,GB to produce the two numbers $d(31,31)=2$, $d(31,32)=38$ that $P_0$
and $P_1$ now give for free---so the pinned polynomials replace exactly the
strips that hardware cannot reach, at the price of the \tierthree{} label.

\section{Companion sequences}
\label{sec:companions}

Counting the same cells under coarser equivalences gives the free, one-sided,
and symmetry-typed companions, all from the fixed count by Burnside's lemma
over the square's symmetry group $D_4$:
\[
  \mathrm{Free}=\tfrac18(\mathrm{Fixed}+2R_{90}+R_{180}+2H+2D),\qquad
  \mathrm{OneSided}=\tfrac14(\mathrm{Fixed}+2R_{90}+R_{180}),
\]
where $R_{90},R_{180},H,D$ count the animals fixed by a $90^\circ$ rotation,
a $180^\circ$ rotation, an axis reflection, and a diagonal reflection
(Method~D). At $n=34$: $R_{90}=0$ (invariance needs $n\equiv0,1\pmod4$;
$34\equiv2$), $R_{180}=\num{48710062997772}$, $H=\num{17385821908346}$. The
derivations assert the $/4$ and $/8$ integrality of the numerators at every
$n$ and validate every published OEIS term of all five companion sequences.

The one-sided extension to $n=34$ is Table~\ref{tab:onesided}. The
free/bilateral/asymmetric companions (A030222, A030234, A030235) and the
free-non-polyomino difference A194596 additionally need $D$, the
diagonal-mirror count, whose exact-bounding-box sweep is the costliest
symmetry (Section~\ref{sec:methodD}): $D$ is banked through $n=28$
($D(28)=\num{47198412763}$). Table~\ref{tab:symcounts} records the
symmetry fixed-point counts themselves, which are of independent interest
(each is an OEIS-eligible sequence).
% TODO(n32-companions): the dmirror n=32 farm completes today: fill the D
% column for n=29..32, then add the companion table
% A030222/A030234/A030235/A194596 for n=25..32 mirroring
% Table~\ref{tab:onesided}. n=33 values, if included, are T3 (P_4-assisted
% strips S>=29) and must carry the label.

\begin{table}[H]
  \centering\small
  \begin{tabular}{r r r r r}
    \toprule
    $n$ & $R_{90}$ & $R_{180}$ & $H$ (axis) & $D$ (diagonal)\\
    \midrule
    25 & \num{6731}   & \num{3161384353}     & \num{2965640141}     & \num{2636626551}\\
    26 & 0            & \num{21807864380}    & \num{8118079636}     & \num{7015813156}\\
    27 & 0            & \num{21570502908}    & \num{20086237186}    & \num{17714707000}\\
    28 & \num{62587}  & \num{149423372035}   & \num{55006721709}    & \num{47198412763}\\
    29 & \num{40607}  & \num{147600215717}   & \num{136445497555}   & ---\\
    30 & 0            & \num{1026290399226}  & \num{373785638266}   & ---\\
    31 & 0            & \num{1012501939560}  & \num{929270191286}   & ---\\
    32 & \num{398124} & \num{7063812264280}  & \num{2546382907164}  & ---\\
    33 & \num{250098} & \num{6960731505456}  & \num{6343327587938}  & ---\\
    34 & 0            & \num{48710062997772} & \num{17385821908346} & ---\\
    \bottomrule
  \end{tabular}
  \caption{Symmetry fixed-point counts, $n=25$--$34$ (\tierone{} to $n=24$
  via the dual-engine gate; \tiertwo{} here). $R_{90}=0$ unless
  $n\equiv0,1\pmod4$. The $D$ column is banked through $n=28$; $n=29$--$32$
  land with the in-flight strip farm.}
  \label{tab:symcounts}
\end{table}

\paragraph{The bilateral identity.}
The bilateral count equals $\tfrac12(H+D)$ exactly, by a double count. Write
the symmetry group of a fixed polyplet as a subgroup $K\le D_4$. If $K$
contains a reflection, its rotations form an index-$2$ subgroup, so exactly
half of $K$'s elements are reflections; otherwise $K$ has none. Count the
incidences $(A,m)$ of a fixed animal $A$ with a reflection
$m\in\operatorname{Stab}(A)$ in two ways. Summed over the four reflections it
is $2H+2D$, since conjugate reflections have equal fixed-counts. Summed over
free classes, a class with stabilizer order $|K|$ contains $8/|K|$ fixed
animals, each with $|K|/2$ reflections in its stabilizer when the class is
bilateral (and none otherwise)---so exactly $4$ incidences per bilateral
class. Hence $2H+2D=4\cdot(\text{bilateral classes})$. This was also checked
numerically against every published term of A030234.

\section{Hole stratification}
\label{sec:holes}

A \emph{hole} of a polyplet is a bounded connected component of its
complement. For king (8-connected) foreground the topologically consistent
choice is \textbf{4-connected} background: a single diagonal staircase of
cells then seals a region (its corner gaps are not 4-passable), matching the
Jordan-curve picture and the convention of
\href{https://oeis.org/A389193}{A389193} for ordinary polyominoes.

We stratify A006770 by hole count: let $T_{\mathrm{hole}}(n,k)$ be the number
of fixed $n$-cell polyplets with exactly $k$ holes, so
$\sum_k T_{\mathrm{hole}}(n,k)=a(n)$. The column transfer matrix carries $k$
as a second additive index via the Euler characteristic
($\#\text{holes}=1-\chi$ for a connected animal, $\chi$ accumulated by local
$2\times2$ bit-quad increments~\cite{gray}); the resulting counts are exact,
byte-identical to an independent per-animal flood fill for $n\le14$, and have
row sums equal to A006770. The triangle is exact through $n=18$; rows
$15$--$18$, beyond the flood oracle, are confirmed by a second-architecture
recount (clang/ARM against GCC/x86, byte-identical), so no row rests on a
single engine. At $n=18$ there are $\num{16503616943998}$ hole-free polyplets
and $\num{4970078092354}$ with a single hole, the maximum being $10$. The
hole-free and one-hole columns are absent from OEIS (checked by Superseeker);
b-files are prepared in the release.

\begin{table}[H]
  \centering
  \resizebox{\textwidth}{!}{%
  \begin{tabular}{r r r r r r r r r r r r}
    \toprule
    $n$ & $k=0$ & $k=1$ & $k=2$ & $k=3$ & $k=4$ & $k=5$ & $k=6$ & $k=7$ & $k=8$ & $k=9$ & $k=10$\\
    \midrule
    1 & \num{1} &  &  &  &  &  &  &  &  &  &  \\
    2 & \num{4} &  &  &  &  &  &  &  &  &  &  \\
    3 & \num{20} &  &  &  &  &  &  &  &  &  &  \\
    4 & \num{109} & \num{1} &  &  &  &  &  &  &  &  &  \\
    5 & \num{622} & \num{16} &  &  &  &  &  &  &  &  &  \\
    6 & \num{3664} & \num{166} & \num{2} &  &  &  &  &  &  &  &  \\
    7 & \num{22094} & \num{1456} & \num{42} &  &  &  &  &  &  &  &  \\
    8 & \num{135609} & \num{11788} & \num{538} & \num{6} &  &  &  &  &  &  &  \\
    9 & \num{843941} & \num{91300} & \num{5596} & \num{144} & \num{1} &  &  &  &  &  &  \\
    10 & \num{5310754} & \num{688034} & \num{52268} & \num{2082} & \num{42} &  &  &  &  &  &  \\
    11 & \num{33724862} & \num{5091820} & \num{457754} & \num{24144} & \num{820} & \num{8} &  &  &  &  &  \\
    12 & \num{215793158} & \num{37209939} & \num{3841776} & \num{248526} & \num{11482} & \num{263} & \num{2} &  &  &  &  \\
    13 & \num{1389673091} & \num{269452496} & \num{31290820} & \num{2374520} & \num{135235} & \num{4808} & \num{96} &  &  &  &  \\
    14 & \num{8998648488} & \num{1937956658} & \num{249306876} & \num{21556520} & \num{1437090} & \num{67020} & \num{2186} & \num{22} &  &  &  \\
    15 & \num{58548155506} & \num{13865448048} & \num{1953275736} & \num{188568632} & \num{14261808} & \num{804084} & \num{35086} & \num{808} & \num{6} &  &  \\
    16 & \num{382526638033} & \num{98796744112} & \num{15103433472} & \num{1603972074} & \num{134769521} & \num{8779911} & \num{465138} & \num{16410} & \num{314} & \num{1} &  \\
    17 & \num{2508473632910} & \num{701664375696} & \num{115556223626} & \num{13348948048} & \num{1227552608} & \num{89826208} & \num{5490242} & \num{251104} & \num{7847} & \num{104} &  \\
    18 & \num{16503616943998} & \num{4970078092354} & \num{876484089730} & \num{109173877126} & \num{10866411352} & \num{876032532} & \num{59950610} & \num{3272674} & \num{137298} & \num{3460} & \num{30} \\
    \bottomrule
  \end{tabular}}
  \caption{The complete hole triangle $T_{\mathrm{hole}}(n,k)$
  ($4$-connected-background convention), exact through $n=18$; row sums
  equal A006770.}
  \label{tab:holes}
\end{table}

\paragraph{Maximum hole area.}
How much empty area can $n$ cells enclose? Let $M(n)$ be the maximum number
of enclosed empty cells over all $n$-cell polyplets. A diamond (L$^1$-ball)
wall of $4r$ cells on $|x|+|y|=r$ is king-connected and seals the $2r^2-2r+1$
cells inside it, so $M(4r)\ge 2r^2-2r+1$ (the centered square numbers,
\href{https://oeis.org/A001844}{A001844}). This diamond is asymptotically
optimal: when the boundary is measured in the king ($L^\infty$, Chebyshev)
metric, the area-maximizing shape is the diamond, the $L^1$ ball dual to the
$L^\infty$ square (Busemann's isoperimetrix~\cite{strang}). Hence
$M(n)\sim n^2/8$, twice the polyomino rate $\sim n^2/16$, since only king
animals can run the diagonal walls the $L^1$ ball requires.

Brute force gives the exact values
$M(1{:}16)=0,0,0,1,1,2,3,5,6,8,10,13,15,18,21,25$, which the diamond attains
at $n=4,8,12,16$. In fact all sixteen values fit a single closed form: they
coincide with \emph{the nearest integer to $(n-2)^2/8$}---the sequence
\href{https://oeis.org/A001971}{A001971} shifted by two---which agrees with
the diamond exactly at every multiple of four, since
$\bigl[(4r-2)^2/8\bigr] = 2r^2-2r+1$.

\begin{conjecture}
\label{conj:diamond}
$M(n)=\left\lfloor \tfrac{(n-2)^2}{8}+\tfrac12 \right\rfloor$ for all
$n\ge1$ (equivalently $M(n)=\mathrm{A001971}(n-2)$), verified here for
$n\le16$. In particular the diamond wall is exactly---not just
asymptotically---optimal at every multiple of four, and away from multiples
of four the optimum tracks the same isoperimetric quota with a two-cell
corner correction.
\end{conjecture}

\begin{openproblem}
\label{op:mn}
Prove Conjecture~\ref{conj:diamond}, or extend the verification past
$n=16$; the first untested value is the prediction $M(17)=28$.
\end{openproblem}

\section{Generating functions by height, and rigorous bounds}
\label{sec:gf}

For a fixed bounding-box height $H$, the strip transfer matrix has a bounded
state space, so $T(n,H)$ satisfies a finite-order linear recurrence and its
generating function $G_H(x)=P_H(x)/Q_H(x)$ is rational~\cite{stanley}. We
recover these exactly for $H\le10$ by Berlekamp--Massey~\cite{massey} over
$\mathbb{Z}/p\mathbb{Z}$ for several primes followed by CRT, each accepted
only when the fitted recurrence reproduces held-out terms. The denominators
$Q_H$ have orders\footnote{These orders are $3$-term running sums of
\emph{atom} degrees $1,2,4,9,29,68,181,\dots$: each $Q_H$ factors as
$N_{H-2}N_{H-1}N_H$, where $N_H$ is the denominator of the height-$H$ strip
generating function, so every irreducible factor divides exactly three
consecutive $Q_H$---a second-difference effect, universal across finite-range
animal classes and verified here for $H\le7$. Neither degree sequence appears
in OEIS (Superseeker-checked).} $1,3,7,15,42,106,278,711,1897,5005$ for
$H=1,\dots,10$; the $H=10$ recovery, needing $\sim\!10{,}000$ series terms,
was reproduced on a second architecture byte-identically. As an example,
\[
  G_3(x)=\frac{9x^3-8x^4-2x^5+4x^6+x^7}{1-7x+15x^2-9x^3-3x^4+5x^5-x^6-x^7}.
\]

\paragraph{Hole-resolved.}
Refining by hole count keeps each fixed-$(H,k)$ row C-finite: we recover
$G_{H,k}$ for $H\le6$ (all reachable $k$), the rows $k\le2$ at $H=7$, and the
row $k=0$ at $H=8$, each validated against a fresh prime. Across the
recovered rows the order of $G_{H,k}$ grows about linearly in $k$ (for
$H=3,\dots,7$ it fits $c_H(k+1)$ with $c_H=6,20,68,185,537$)---recorded as an
empirical regularity. Summed over all heights the columns, like $a(x)$
itself, are not expected to be D-finite, so the rationality is genuinely a
fixed-height phenomenon.

\paragraph{Rigorous lower bounds.}
Summing the fixed-height GFs gives $a(n)\ge\sum_{H\le 10}T(n,H)$ exactly for
any $n$, by extracting coefficients from rational functions. When this bound
was computed, $a(25)$ was out of reach and the bound gave
$a(25)\ge\num{6657105903966723073}$; the exact value
$\num{14994811325186658577}$, computed later by Method~C, confirms it (the
bound captured $44\%$)---a retrospective certification of the GF machinery.
The bound remains available for any $n$ beyond the current frontier, though
its capture fraction shrinks as taller strips dominate.
% TODO(gf-bound-n35plus): recompute the showcase bound at n=35 or 40 via
% paper/gf_bound.py and quote the capture fraction against a(34)'s trend.

\section{Sampling portrait}
\label{sec:sampling}

The counting DP doubles as an exact uniform sampler~\cite{nijenhuis}: a
single forward pass through the transfer matrix, selecting each successive
column with probability proportional to the number of completions it admits,
draws an animal uniformly at random. (The portrait below was made at the
$a(19)$/$a(20)$ stage with Method B, whose per-state completion counts the
sampler needs; the production spill engine does not retain them, so the
portrait stays at that scale.)

We validated the sampler on $a(19)$: an independent king-move BFS confirmed
every specimen was connected, and the sampled distribution of heights
reproduced $T(19,H)/a(19)$. Ten thousand samples against a state space of
$1.5\times10^{14}$ can rule out only gross bias; no goodness-of-fit power is
claimed against subtle non-uniformity in the transition probabilities, and
the portrait below should be read at that resolution---qualitative
character, not calibrated statistics. A $10{,}000$-specimen sample shows the typical
$19$-cell polyplet is \emph{sparse, branchy, corner-glued}: about ten
rook-connected pieces held together mostly by corners ($59\%$ of contacts),
$\sim\!25\%$ box density, and a $28\%$ chance of enclosing a hole. It is
essentially never a polyomino: the exact frequency
$A001168(19)/a(19)\approx4\times10^{-5}$ predicts about $0.4$ such specimens
per $10{,}000$, and none appeared. Sampling at several sizes
($n=8,\dots,19$) and fitting the radius of gyration gives $R_g\sim n^{\nu}$
with $\nu\approx0.68$, above the accepted two-dimensional value
$\nu\approx0.6408$~\cite{jensen2001} as a fit over so short a range should
be; the holed fraction rises from $8.3\%$ to $28.3\%$ over $n=8\to19$.
These are heuristic observations, not estimates with error control.

\section{Reach and limits}
\label{sec:reach}

Method C's cost structure is set by its tallest surviving sweep height: at
fixed $H$ the strip's state space is bounded and cost grows only about
linearly in $n$, while adding a unit of height multiplies it by roughly the
per-height frontier growth. With the diagonals $k\le15$ in closed form, the
tallest real sweep stayed at $H=18$ from $a(33)$ to $a(34)$, and the banked
totals grew gently ($22$ to $59$ core-hours over $a(30)$--$a(34)$; see
Appendix~\ref{sec:profile}). $a(35)$ is therefore within easy hardware
reach---an $H=19$ sweep at roughly $3\times$ the $H=18$ cost---but faces a
verification gap rather than a compute gap: either sweep $H=19$ for real and
check it against nothing, or wire $P_{16}$, which cannot be independently
held out until $a(35)$ itself exists. The symmetric counts scale as
$\sqrt{a}$-class problems and reached $n=34$ for three of the four types;
the diagonal-mirror type, whose forced-square bounding box compresses least,
is the binding constraint on the remaining companions---its per-strip RAM
peaked at $83$\,GB at $n=32$ after a compact-frontier rewrite (flat sharded
arrays, ranged count vectors, $\approx1.6\times$ RAM and $\approx2\times$
speed over the hash-map frontier it replaced), and the pinned
quasi-polynomials of Section~\ref{sec:dmdiag} substitute for the strips
beyond hardware at $n\ge33$ under the \tierthree{} label.

\section{Data and code availability}
\label{sec:availability}

The source (all engines and gates), the per-height triangle rows for every
term through $n=34$ with per-term provenance ledgers, the b-files for
A006770 and the companion sequences, the hole triangle and its columns, the
recovered generating functions, and the sample statistics are in the
accompanying repository; derived artifacts are regenerated on demand by the
included tools. The cross-checks are encoded as runnable gates, and a
self-contained checker re-verifies the paper's numerical claims against the
committed and regenerated data.
% TODO(verify-claims): paper/verify_claims.py still targets the a(19)/a(20)
% claim set; re-point it at this paper's tables (a(n) to 34, A030233 to 34,
% the n=32 companions when banked, P_k tables) before submission.
% TODO(bfile-34): results/b006770_upload.txt currently ends at n=32; extend
% to 34 from results/ns_a34/triangle.txt.

\section{Author's note: use of AI}
\label{sec:ai}

This is an amateur contribution by an unaffiliated author in collaboration
with an AI assistant (Anthropic's Claude Code, using the Claude Opus~4.8 and
Claude Fable~5 models over the project's lifetime), which helped with
algorithm design, implementation, running the computations, and drafting.
The author set the goals, made the design and verification decisions, owns
and operated the machines, inspected the results, and is responsible for
their correctness.

The verification protocol of Section~\ref{sec:confirm}---independent
algorithms where they exist, decorrelated recounts and held-out closed-form
checks where they do not, and explicit confidence tiers throughout---ensures
the central claims do not rest on the correctness of any single
\emph{program}. They do rest, past $n=19$, on a single algorithm family;
the tier system exists so that dependence is never hidden. All code and data
are released for independent checking.

\appendix
\section{Computational profile}
\label{sec:profile}

For future benchmarking, the measured costs of the headline computations.
All machines are commodity hardware owned and operated by the author:
\emph{gympie} (Apple silicon, 10 performance cores, 24\,GB), \emph{ayr}
(AMD x86-64, 32 cores, 78\,GB), \emph{dalby} (Arm Neoverse-N1, 80 cores,
128\,GB).

\begin{table}[H]
  \centering\small
  \begin{tabular}{l r r r}
    \toprule
    term (full triangle, Method C, dalby) & wall (h) & core-hours & peak RSS (GB)\\
    \midrule
    $a(30)$ & 1.0 & 22.1 & 1.7\\
    $a(31)$ & 2.8 & 44.8 & 5.8\\
    $a(32)$ & 3.0 & 49.0 & 5.6\\
    $a(33)$ & 3.4 & 53.5 & 5.9\\
    $a(34)$ & 3.7 & 58.5 & 6.3\\
    \bottomrule
  \end{tabular}
  \caption{Banked per-term cost profiles for the uniform final-engine runs
  (earlier terms were computed by intermediate engine generations and are
  not directly comparable). The gentle growth reflects the closed-form
  diagonals holding the tallest real sweep at $H=18$ throughout this range;
  RAM stays small because the frontier spills to disk.}
  \label{tab:costs}
\end{table}

\begin{itemize}
  \item \textbf{$a(34)$, full triangle (Method C, dalby):} wall
        $\approx3.7$\,h across all 34 heights; $\approx59$ core-hours; peak
        resident memory $6.4$\,GB (the frontier spills to disk, so RAM
        stays small; the tallest surviving real sweep height dominates,
        $\sim70\%$ of the wall). The cross-ISA recount on ayr reproduced
        the triangle byte-identically.
  \item \textbf{Diagonal-mirror symmetry count, $n=32$ (Method D, farmed):}
        the six deepest exact-bounding-box strips ($S=26$--$31$, dalby)
        cost $3.0$\,M CPU-seconds $\approx835$ core-hours with per-strip
        peak memory up to $82.7$\,GB---an order of magnitude more compute
        than $a(34)$ itself, which is why the quasi-polynomials of
        Section~\ref{sec:dmdiag} matter: they replace exactly the strips
        whose memory grows fastest.
        % TODO(n32-cost): add the ayr tail (S<=25) totals when the farm
        % completes, and the n=33 run's totals when banked.
  \item \textbf{Symmetric counts at $n=34$ (Method D):} axis mirror
        $\approx3.3$\,h wall / $53$\,k CPU-seconds / $8.2$\,GB;
        $180^\circ$ rotation $\approx51$\,min wall / $18$\,k CPU-seconds /
        $8.5$\,GB (both dalby); the $90^\circ$ count runs in seconds on
        gympie.
\end{itemize}

\begin{thebibliography}{9}

\bibitem{redelmeier}
D.~H. Redelmeier,
\emph{Counting polyominoes: yet another attack},
Discrete Mathematics \textbf{36}(2) (1981), 191--203.

\bibitem{klarner}
D.~A. Klarner,
\emph{Cell growth problems},
Canadian Journal of Mathematics \textbf{19} (1967), 851--863.

\bibitem{jensen2001}
I.~Jensen,
\emph{Enumerations of lattice animals and trees},
Journal of Statistical Physics \textbf{102}(3--4) (2001), 865--881.

\bibitem{jensen2003}
I.~Jensen,
\emph{Counting polyominoes: a parallel implementation for cluster computing},
in \emph{Computational Science --- ICCS 2003}, LNCS \textbf{2659}, Springer,
2003, pp.~203--212.

\bibitem{bousquet}
M.~Bousquet-M\'elou,
\emph{A method for the enumeration of various classes of column-convex polygons},
Discrete Mathematics \textbf{154}(1--3) (1996), 1--25.

\bibitem{nijenhuis}
A.~Nijenhuis and H.~S. Wilf,
\emph{Combinatorial Algorithms for Computers and Calculators}, 2nd~ed.,
Academic Press, 1978, pp.~99--116.

\bibitem{stanley}
R.~P. Stanley,
\emph{Enumerative Combinatorics, Volume~1}, 2nd~ed.,
Cambridge University Press, 2012 (rational generating functions,
quasi-polynomials, and the transfer-matrix method, Ch.~4).

\bibitem{massey}
J.~L. Massey,
\emph{Shift-register synthesis and BCH decoding},
IEEE Transactions on Information Theory \textbf{15}(1) (1969), 122--127.

\bibitem{gray}
S.~B. Gray,
\emph{Local properties of binary images in two dimensions},
IEEE Transactions on Computers \textbf{C-20}(5) (1971), 551--561.

\bibitem{barequet}
G.~Barequet, G.~Rote, and M.~Shalah,
\emph{$\lambda > 4$: an improved lower bound on the growth constant of polyominoes},
Communications of the ACM \textbf{59}(7) (2016), 88--95.

\bibitem{strang}
G.~Strang,
\emph{Maximum area with Minkowski measures of perimeter},
Proceedings of the Royal Society of Edinburgh Section A: Mathematics
\textbf{138}(1) (2008), 189--199.

\bibitem{oeis}
OEIS Foundation Inc.,
\emph{The On-Line Encyclopedia of Integer Sequences} (2026),
sequences A006770, A030222, A030233, A030234, A030235, A194596, A001168,
A389193, A001844.
\url{https://oeis.org}.

\end{thebibliography}

\end{document}
